\section*{الفصل \ref{ch:g2:drawing-to-map} --- من الرسم إلى الخريطة}

\begin{solution}{exo:g2:drawing-to-map:1}
كما يُرى من الأعلى.
\end{solution}

\begin{solution}{exo:g2:drawing-to-map:2}
الخطوط الزرقاء على الجدار الغربي هي النوافذ. وفي الركن الشمالي الشرقي
الخزانة.
\end{solution}

\begin{solution}{exo:g2:drawing-to-map:3}
عنوان، \omterm{def:g2:drawing-to-map:legend}{ومفتاح}، \omterm{def:g2:drawing-to-map:scale}{ومقياس رسم}، \omterm{def:g2:drawing-to-map:directions}{ووردة رياح} (أو سهم يشير إلى الشمال).
\end{solution}

\begin{solution}{exo:g2:drawing-to-map:4}
الشمال في الأعلى؛ والغرب على اليسار.
\end{solution}

\begin{solution}{exo:g2:drawing-to-map:5}
المربع الواحد يمثل مترًا واحدًا: طول الطاولة متران وعرضها متر واحد.
\end{solution}

\begin{solution}{exo:g2:drawing-to-map:6}
المربع البني هو المزرعة؛ والشكل الأخضر هو الحرج.
\end{solution}

\begin{solution}{exo:g2:drawing-to-map:7}
الحرج جنوب \omterm{def:g1:landscapes-around-us:water}{النهر}، وشرق المزرعة.
\end{solution}

\begin{solution}{exo:g2:drawing-to-map:8}
المربع الواحد يمثل كيلومترًا واحدًا: 10 كيلومترات من الغرب إلى الشرق و8
كيلومترات من الجنوب إلى الشمال.
\end{solution}

\begin{solution}{exo:g2:drawing-to-map:9}
ثلاثة صفوف في كل منها ثلاث طاولات: 9 طاولات. والباب في الجزء الجنوبي من
الجدار الشرقي، إذن الطاولة الأقرب إلى الباب في الجنوب (الركن الجنوبي
الشرقي من الغرفة).
\end{solution}

\begin{solution}{exo:g2:drawing-to-map:10}
طول السرير مربعان. والغرفة 4 مربعات في 3 مربعات: 3 صفوف في كل منها 4
مربعات، أي $4 + 4 + 4 = 12$ مربعًا.
\end{solution}

\begin{solution}{exo:g2:drawing-to-map:11}
«اليسار» يتوقف على الجهة التي تتجه إليها لينا: فإن استدارت صار \omterm{def:g1:landscapes-around-us:water}{النهر} على
يمينها. أما عند الجميع: «\omterm{def:g1:landscapes-around-us:water}{النهر} شمال المزرعة».
\end{solution}
